%%%%%%%%%%%%%%%%%%%%%---Chapter7 Formatting---%%%%%%%%%%%%%%%%%%%%%%%%%%%%%%%%%%
\documentclass[Thesis.tex]{subfiles}
\begin{document}
\setcounter{chapter}{6} %%Chapter Numbering


%%%%%%%%%%%%%%%%%%%%%%%%%%%%%%%%%%%%%%%%%%%%%%%%%%%%%%%%%%%%%%%%%%%%%%%%%%%%

\chapter{Conclusions and Future Scope}\label{chap7:conclusion}

\section{Conclusions}

This thesis evaluated hybrid quantum--classical models for quantum error
mitigation. The benchmark compares ten main expectation-value mitigation methods
on $4500$ reference circuit--observable pairs, under an idealised Markovian noise
model and a device-informed non-Markovian one. In these experiments, no hybrid
model showed a consistent accuracy advantage over the strongest comparison
methods.

\subsection{The matched kernel comparison}

On random held-out circuits in the idealised regime, classical RBF reaches
$0.0404$, the fidelity kernel reaches $0.0421$, and
linear zero-noise extrapolation reaches $0.0737$. The paired difference
RMSE(RBF) $-$ RMSE(fidelity) is $-0.00168$ (95\% CI $[-0.00324,\allowbreak -0.00071]$). Across ten repeated
grouped holdouts, the mean and standard deviation are $0.0416\pm0.0011$
for RBF and $0.0426\pm0.0012$ for fidelity. The per-split paired intervals
favour RBF in $9$ splits and favour fidelity in $0$; the tenth interval includes zero.
At the selected bandwidth the fidelity Gram matrix has correlation $0.993$ with a
Gaussian kernel on the same features, and the RBF grid contains a Gaussian of almost
the same width. The matched RBF control can therefore represent nearly the same
similarity structure, which is consistent with RBF matching or beating the fidelity
kernel in these experiments. With $31$ reference circuits the fidelity kernel is
favoured in $9$ of $20$ draws and RBF in $4$; at $63$ circuits the counts are $5$
and $8$.

\subsection{Sensitivity to hyperparameters and initialisation}

The fixed-bandwidth diagnostic gives a fidelity-kernel RMSE of
$0.1526$ at $\beta=1$ and $0.0421$ at the best observed
bandwidth. At $\beta=1$ the mean off-diagonal kernel entry is $0.204$, far above
the concentrated limit, so the untuned failure comes from a narrow kernel. The separate
variational seed study gives $0.0935 \pm 0.0327$ on Split~A for the data-reuploading
VQC, with five of ten seeds worse than no mitigation, and the single main-table fit is
not reproducible across scripts. Both results motivate systematic tuning and repeated
initialisations in benchmarks.

\subsection{The two noise-scaling protocols}

In the idealised noise-strength sweep, RBF's reduction relative to
unmitigated RMSE narrows from $59.6\%$ at $0.5\times$ to $24.6\%$ at $4\times$
on the same held-out circuits.

On the realistic reference pool, RBF and fidelity reach $0.0622$ and
$0.0622$ under parameter scaling, and $0.0762$ and
$0.0757$ under folding on the random holdout. Both paired
RBF-minus-fidelity intervals include zero.
On that holdout, the kernels' RMSE gap below linear ZNE
narrows from about $0.040$ with parameter scaling to $0.014$ with folding.
On the depth split, the paired comparison favours RBF under parameter
scaling. The folding interval is
$[-0.00158,\allowbreak +0.00289]$.
Both protocols share their native-scale shots, and folded ZNE uses the
actual ratios of executed to original layers.

In the three-pool replication, folding lowers linear ZNE error and raises
both fitted kernels' RMSE relative to parameter scaling on both splits in
all nine shot-draw combinations (Section~\ref{sec:ch5_folding_seeds}).
The reset intervention and drift comparisons give further controls for the
same simulated device.

\subsection{Finite-shot kernels and circuit cost}

On the parameter-scaled random holdout, PSD repair on each training block defines
the projection for its validation or test cross-kernel. Mean RMSE is $0.0658$
at $10^5$ shots per entry and $0.0629$ with exact entries, against $0.0622$ for
RBF (Section~\ref{sec:ch5_shotkernel}). Both fidelity variants use the same
folds, bandwidth and ridge grid. At $10^5$ shots per entry, tuning, fitting and
scoring the fidelity kernel on this split use $2.78\times10^{12}$ simulated
fidelity executions, and each subsequent prediction adds $3.15\times10^{8}$.
RBF computes its similarities classically from the $12{,}288$ input executions
per row that both kernels share.

\subsection{Shot noise and analytic extrapolation}

On idealised random circuits, the kernels reach $0.0029$ and $0.0031$ RMSE with
shot-free inputs against $0.0411$ and $0.0427$ at $4096$ shots, so most of their
error at the default budget comes from input shot noise. Exponential extrapolation
reaches $0.0119$ with shot-free inputs but $0.1548$ at $4096$ shots, and it returns
the linear value on $317$ of the $1350$ Split~A test rows. In the device-informed
regime the residual error of linear extrapolation is mainly bias: on shot-free
production values it reaches $0.0997$ under parameter scaling against $0.1206$
unmitigated.

\subsection{Non-Markovian memory and calibration drift}

On the realistic random holdout, removing the defect lowers the unmitigated RMSE from
$0.1239$ to $0.0807$, while resetting it after every layer leaves $0.1247$. Most of the
defect's damage to the raw measurement therefore comes from energy exchange with the
defect rather than from the state it retains. The retained state does affect the
learned mitigators: the production-minus-reset RMSE difference is $+0.0095$ for RBF
and $+0.0098$ for fidelity under parameter scaling, and $+0.0211$ and $+0.0204$ under
folding, all with paired intervals excluding zero. Under parameter scaling the scaled
settings also change the coherent dynamics, so the folded contrast is the more direct
measure of memory.

After the simulated $30\%$ recalibration, the unmitigated RMSE on the same held-out
circuits rises from $0.1239$ to $0.1298$. Ridge regression ($+0.0034$) and linear
extrapolation ($+0.0043$) show resolved increases; the other learned arms show none.
This drift changes the raw error by about $5\%$, so it tests a mild recalibration.

\subsection{Feature support and depth transfer}

On the idealised deeper-circuit split, RBF reaches $0.0658$
and fidelity reaches $0.1433$, using the same six inputs. The raw measurement
reaches $0.1217$. Removing depth and two-qubit gate count while retaining
observable weight gives $0.0718$ for RBF and $0.0723$ for fidelity. This feature
ablation changes the inputs and fidelity-map width, and its result applies to
the tested depth split (Section~\ref{sec:ch5_robustness}).

Under these shifts the model classes separate. With all six inputs on deeper circuits,
RBF, the boosted trees and ridge regression stay $41$--$46\%$ below the unmitigated
error, while the neural network, the fidelity kernel and both variational circuits are
worse than no mitigation. Ridge regression is the lowest-error arm from
hardware-efficient to Ising circuits ($0.0546$), where RBF and fidelity are $2.0\%$ and
$19.8\%$ worse than no mitigation. The arms that degrade gracefully in these tests are
ridge regression, which extrapolates linearly, and the boosted trees, which hold their
boundary values outside the training range; RBF does so on the depth split but not
from hardware-efficient to Ising circuits.

\subsection{Scope and limitations}

Mitigation uses simulated circuits with at most six qubits. The device-informed
regime models one set of noise parameters and one drift condition, including a
coupled defect on qubit~$0$, which appears in four of the five measured Pauli
strings. The model omits readout assignment errors and leakage. Parameter scaling in this
regime also scales the coherent over-rotation, crosstalk and defect coupling, and the
fidelity kernel uses six qubits at every system size. The trace-distance witness covers two
probe pairs on an idle register. The reset comparison studies a specific
between-layer intervention at fixed collision and damping parameters.
The fidelity resource estimates count simulated circuit executions with ideal
overlap probabilities and binomial sampling. The analytic fidelity reference
computes exact state overlaps in simulation.

The initial sweeps and cross-family experiment use one data split per setting.
The ten repeated holdouts are drawn from one dataset, so their test sets overlap and are not independent. The pool replication uses three circuit
pools, and shot draws within a pool share its circuits. Intervals across shot
budgets and models are pointwise.

\section{Future scope}

\begin{enumerate}
    \item \textbf{Hardware data through the same pipeline.} Use measured
    expectations with a simulable or Clifford-replacement reference. Keep
    all observables measured on a circuit in the same group and build inputs
    as in Section~\ref{sec:ch3_features}. Use matched native measurements
    and executed folding ratios to compare learned models with extrapolation
    on a processor.

    \item \textbf{Support-aware mitigation.} The depth ablation motivates
    explicit checks for coverage of structural inputs. Compare an input-support
    check, a shape-constrained regressor and a fallback to analytic extrapolation
    on depth or family shifts.

    \item \textbf{Hardware training of the variational arms.} The variational
    arms use simulator backpropagation. For the standard two-shift identity, a
    hardware gradient uses two shifted settings per differentiated gate angle
    and circuit input, before accounting for measurement shots. Compare a
    shot-resolved training and prediction budget with linear extrapolation.

    \item \textbf{Wider noise and size sweeps.} Repeat the paired noise-strength
    sweep on additional circuit pools and extend the system-size study to
    test how both kernels scale. Beyond density-matrix simulation, these
    experiments require controlled approximations or hardware measurements.

\end{enumerate}

\section{Closing remark}

In these simulations, a matched classical RBF kernel reaches or beats the fidelity
kernel, and at its selected bandwidth the fidelity kernel is close to a Gaussian kernel.
The ranking of all methods depends on the circuit distribution, the feature set and the
noise-scaling protocol. A hardware deployment would also have to pay for the reference
data, the input measurements and the fidelity-estimation circuits counted here.

\end{document}
